\section{The Chart Family: A Grammar-of-Graphics Layer}
\label{sec:chart-family}

The chart family is the one genuinely new architectural component required by the
Mermaid-superset strategy. Unlike node-link, timeline, and tree families (which map
\emph{structural} entities to \emph{positioned shapes}), charts map \textbf{quantitative
data} to \textbf{geometric marks} through scales and coordinate systems.

This section specifies the chart grammar drawing on the Grammar of
Graphics~\cite{grammarofgraphics2005} and Vega-Lite~\cite{vegalite2017} traditions,
adapted to our two-IR-layer architecture.

\subsection{Design Principles for Charts}

\begin{enumerate}
  \item \textbf{Data $\to$ geometry, not data $\to$ pixels.} The Domain IR specifies data
        and encoding; the layout engine computes geometry; the theme styles it.
  \item \textbf{Mermaid-compatible syntax.} Each chart type's Mermaid syntax parses
        faithfully to our Chart IR.
  \item \textbf{Deterministic layout.} All scale computations, tick placement, and mark
        positioning are closed-form (no iteration, no randomness).
  \item \textbf{Theme-driven aesthetics.} Colors, fonts, gridlines, and labels are
        theme-controlled. The IR carries data and encoding, never style.
\end{enumerate}

\subsection{Chart Domain IR}

\begin{lstlisting}[language=javascript,caption={Chart Domain IR shape (TypeScript)}]
interface ChartDocument {
  version: string;
  chartType: 'pie' | 'xy' | 'quadrant' | 'radar';
  title?: string;
  data: ChartData;
  encoding: ChartEncoding;
  config?: ChartConfig;
}

interface ChartData {
  values: Array<Record<string, number | string>>;
}

interface ChartEncoding {
  // Type-specific encoding channels
  x?: FieldEncoding;      // xy charts
  y?: FieldEncoding;      // xy charts
  color?: FieldEncoding;  // categorical color
  size?: FieldEncoding;   // mark size (radar axis values)
  theta?: FieldEncoding;  // pie angle
  label?: FieldEncoding;  // text labels
}

interface FieldEncoding {
  field: string;          // key into data values
  type: 'quantitative' | 'nominal' | 'ordinal';
  scale?: ScaleConfig;
}
\end{lstlisting}

\subsection{Chart Types}

\subsubsection{Pie Chart}

Maps categorical values to angular sectors. Mermaid syntax:

\begin{Verbatim}[fontsize=\small]
pie title Pet Adoption
    "Dogs" : 386
    "Cats" : 85
    "Rats" : 15
\end{Verbatim}

\noindent Domain IR encoding: \texttt{theta} channel maps value to sector angle;
\texttt{color} channel maps label to fill color (theme palette).

\subsubsection{XY Chart (Bar/Line)}

Cartesian charts with categorical or quantitative axes. Mermaid syntax:

\begin{Verbatim}[fontsize=\small]
xychart-beta
    title "Monthly Revenue"
    x-axis [Jan, Feb, Mar, Apr, May]
    y-axis "Revenue (USD)" 0 --> 5000
    bar [1200, 2400, 3100, 2800, 4200]
    line [1000, 2200, 2900, 2600, 4000]
\end{Verbatim}

\noindent Domain IR: \texttt{x} encodes categories or values; \texttt{y} encodes quantitative
measure; mark type (bar/line) stored in encoding config.

\subsubsection{Quadrant Chart}

Four-quadrant scatter plot with labeled axes and positioned items. Mermaid syntax:

\begin{Verbatim}[fontsize=\small]
quadrantChart
    title Reach and engagement
    x-axis Low Reach --> High Reach
    y-axis Low Engagement --> High Engagement
    quadrant-1 We should expand
    quadrant-2 Need to promote
    Campaign A: [0.3, 0.6]
    Campaign B: [0.45, 0.23]
\end{Verbatim}

\noindent Domain IR: \texttt{x} and \texttt{y} encode position on [0,1] scales;
quadrant labels are metadata.

\subsubsection{Radar Chart}

Multi-axis radial chart. Mermaid syntax (proposed):

\begin{Verbatim}[fontsize=\small]
radar
    title Skills Assessment
    axes: [Speed, Reliability, Comfort, Safety, Efficiency]
    "Model A": [4, 3, 2, 5, 4]
    "Model B": [2, 5, 4, 3, 3]
\end{Verbatim}

\noindent Domain IR: axes define radial dimensions; each series maps values to axis positions.

\subsection{Layout Engine}

The chart layout engine is shared across all four chart types:

\begin{enumerate}
  \item \textbf{Scale computation.} From data domain (min/max or categories) to pixel range,
        using linear, band, or radial scales. All scales are closed-form.
  \item \textbf{Axis/grid generation.} Tick positions, labels, and gridlines computed
        deterministically from scale + theme config.
  \item \textbf{Mark placement.} Data points mapped through scales to Scene IR primitives:
        Rect (bars), Path (lines, pie sectors, radar polygons), Circle (scatter points).
  \item \textbf{Label placement.} Deterministic label positioning with collision avoidance
        (priority-based, not iterative).
\end{enumerate}

\subsection{Relation to Vega-Lite}

The chart grammar draws on Vega-Lite's decomposition (data, transform, mark, encoding,
scale, axis, legend) but is deliberately simpler:

\begin{itemize}
  \item No data transforms (filter, aggregate, window)---data arrives pre-aggregated.
  \item No interaction or selection.
  \item No composition operators (layer/concat/facet)---composition is handled by the
        poster layer (Part~\ref{part:composition}).
  \item Subset of mark types: rect, line, arc, point, text.
\end{itemize}

This simplification keeps the chart Domain IR small enough for reliable agent generation
while covering the four Mermaid chart types completely.

\subsection{Implementation Status}

\textbf{Not yet built.} The chart family is the Tier-2 deliverable. It requires:
\begin{itemize}
  \item New Chart Domain IR and JSON Schema
  \item Scale computation library (linear, band, radial)
  \item Chart layout engine compiling to Scene IR
  \item Mermaid parsers for pie, xychart, quadrant, radar syntax
  \item Chart-specific theme extensions (axis styling, gridlines, palettes)
\end{itemize}

The existing Scene IR and backends (SVG, PNG) require no changes---charts compile to the
same Rect/Path/Text/Circle primitives used by all other families.
